%! @@TITLE@@ — Unit @@UNITINT@@, Lesson @@LESSONID@@ — Slides (Beamer)
\documentclass[aspectratio=169, 11pt]{beamer}
\usepackage{@@PREFIX@@-beamer}   % provides \ceruleanheader{} and \sectionlabel[color]{}

\begin{document}

% TITLE SLIDE (CourseName is not defined in beamer, so the name is literal)
{
\setbeamercolor{background canvas}{bg=cerulean}
\begin{frame}[plain]
  \vspace{2.8cm}\hspace{0.55cm}%
  \begin{minipage}{0.9\textwidth}
    {\color{goldacc}\bfseries\small LESSON @@LESSONID@@}\\[0.5cm]
    {\color{white}\bfseries\LARGE @@TITLE@@}\\[1.2cm]
    {\color{frostmid}\small @@COURSENAME@@~$\cdot$~Unit @@UNITINT@@}
  \end{minipage}
\end{frame}
}

% The deck follows the GRADUAL-RELEASE flow (~12 frames):
%   learning targets → warm-up → hook → one frame per numbered notes section (I DO)
%   → We Do → independent-practice launch → debrief → close & assign.
% THERE IS NO SPOILER RULE: name the formal vocabulary on every frame it belongs on.
% \redink{} is for the term the moment it is first attached to an idea.
\newcommand{\redink}[1]{{\color{redacc}\bfseries #1}}

% LEARNING TARGETS — name the vocabulary outright, worded as on the cover.
\begin{frame}[t, plain]
  \ceruleanheader{Today: TODO topic}
  \sectionlabel{Learning targets}
  By the end of today, I can\ldots
  \begin{itemize}\setlength{\itemsep}{6pt}
    \item TODO: target 1, naming the formal term.
  \end{itemize}
  \begin{block}{How today works}
    Warm-up $\to$ \textbf{I work one} $\to$ \textbf{we work one} $\to$
    \textbf{you work TODO on your own} $\to$ we go over them.
  \end{block}
\end{frame}

% WARM-UP
\begin{frame}[t, plain]
  \ceruleanheader{Warm-Up: things you already know}
  \sectionlabel{5 minutes --- alone, then check with a neighbor}
  % TODO: the warm-up items, and which notes section picks each one up again.
\end{frame}

% HOOK — leave it UNRESOLVED if the lesson's crux is the misconception it plants.
\begin{frame}[t, plain]
  \ceruleanheader{TODO: the hook question}
  \sectionlabel[goldacc]{Commit to an answer --- we will settle this later}
  % TODO: the claim or question, and the vote/commitment to take. Do not resolve it here.
\end{frame}

% NOTES SECTION — I DO. Repeat this frame once per numbered notes section (4-6 of them).
% Flag the CRUX section's frame with \sectionlabel[redacc]{}.
% These frames carry the ANSWERS the class writes down — they are the board, not a preview,
% so advance them in step with the handout.
\begin{frame}[t, plain]
  \ceruleanheader{TODO: notes section N title}
  \sectionlabel{I do --- copy this down}
  % TODO: the worked example exactly as it goes on the board, with the formal term in
  % \redink{} the first time it is attached.
\end{frame}

% WE DO — a LIVE surface. Show the task and the data; leave the reasoning BLANK and fill it
% in from what the class supplies. Do NOT ship this frame pre-answered.
\begin{frame}[t, plain]
  \ceruleanheader{We Do: TODO task}
  \sectionlabel[goldacc]{You tell me every step --- I hold the pen}
  % TODO: the task and its data only. Same task the practice box will release, new numbers.
\end{frame}

% INDEPENDENT-PRACTICE LAUNCH — the item range and the rule ONLY.
% DO NOT put the practice items or their answers on a slide before the practice; a deck that
% shows the answer pattern turns the release into copying.
\begin{frame}[t, plain]
  \ceruleanheader{Your turn}
  \sectionlabel[goldacc]{Practice box, items 1--TODO $\cdot$ 10 minutes}
  \begin{block}{The rule}
    \textbf{Silently, on your own.} Every answer ends in a sentence that says what it means
    \emph{in this situation}.
  \end{block}
\end{frame}

% DEBRIEF — walk the displayed practice work, then the cold check.
\begin{frame}[t, plain]
  \ceruleanheader{Going over it}
  \sectionlabel[redacc]{Correct your own work in a second color}
  % TODO: the displayed items in order, the crux item with two disagreeing answers, then the
  % whole-class check-for-understanding question asked cold. This replaces the exit ticket.
\end{frame}

% CLOSE & ASSIGN
\begin{frame}[t, plain]
  \ceruleanheader{What changed today \& what's next}
  \sectionlabel[goldacc]{Homework --- this one is scored}
  % TODO: one line on what changed today. Then name the assignment and say plainly which form
  % it takes — the packet's homework pages OR a specific DESMOS activity (name it). Give the
  % due date and how it is scored; close with a one-line preview of the next lesson.
\end{frame}

\end{document}
